% !TEX root = ../../../main.tex
\section{Sequences of functions}
\label{sec:analysis-i:sequences-of-functions}

% BEGIN TUFTE PLACEHOLDER: supplied sample, lines 1034--1050.
\begingroup\sloppy
% Replace this entire specimen block with reviewed course notes.
\label{template:sec:marginal-material}

Marginal material includes sidenotes, citations, margin notes, and captions.
Normally, the justification of the marginal material follows the justification
of the body text.  If you specify the \docclsopt{justified} document class
option, all of the margin material will be fully justified as well.  If you
don't specify the \docclsopt{justified} option, then the marginal material will
be set ragged right.

You can set the justification of the marginal material separately from the body
text using the following document class options: \docclsopt{sidenote},
\docclsopt{marginnote}, \docclsopt{caption}, \docclsopt{citation}, and
\docclsopt{marginals}.  Each option refers to its obviously corresponding
marginal material type.  The \docclsopt{marginals} option simultaneously sets
the justification on all four marginal material types.

\lipsum[1]
\par\endgroup
% END TUFTE PLACEHOLDER
